

\chapter{Cod HTML (Structura interfeței)}
\label{anexa:html}

Această anexă reunește fragmentele reprezentative din pagina principală
\texttt{index.html}, scrise de absolvent. Pentru concizie, stilurile inline
extinse au fost omise (acestea se află în fișierul \texttt{css/style.css}),
păstrând structura elementelor și identificatorii la care fac referire
modulele JavaScript din Anexa~B.

\section*{Structura paginii principale}

\begin{lstlisting}[language=HTML, caption={Scheletul interfetei: panou lateral, harta si controale (extras din index.html)}, label={lst:htmlskeleton}]
<!DOCTYPE html>
<html lang="en">
<head>
    <meta charset="UTF-8">
    <meta name="viewport"
          content="width=device-width, initial-scale=1.0, viewport-fit=cover">
    <title>ETTIway - Campus Navigation</title>
    <link rel="stylesheet" href="css/style.css">
</head>
<body>
    <button id="sidebar-toggle" class="sidebar-toggle">&#9776;</button>
    <div id="sidebar-overlay" class="sidebar-overlay"></div>

    <!-- Cautare mobila: comutator Cladire / Camera -->
    <div id="mobile-search-container" class="mobile-search-container">
        <div class="segmented-control">
            <button class="segment active" id="segment-building">Cladire</button>
            <button class="segment" id="segment-room">Camera</button>
        </div>
        <div class="search-container mobile-search-box">
            <input type="text" id="mobile-single-search"
                   placeholder="Cauta cladire pe ETTIway" />
        </div>
    </div>

    <!-- Panou lateral (devine bottom sheet pe mobil) -->
    <div id="sidebar">
        <div class="sidebar-header">
            <img src="Icons/logo.svg" alt="ETTIway Logo" class="sidebar-logo">
            <p class="subtitle">Campus Navigation</p>
            <p id="logged-user" class="logged-user" style="display: none;"></p>
            <form action="/logout" method="post" class="logout-form">
                <button type="submit" class="logout-btn">Logout</button>
            </form>
        </div>
        <div class="search-container desktop-search-container">
            <div class="segmented-control desktop-segmented-control">
                <button class="segment active" id="segment-building-desktop">Cladire</button>
                <button class="segment" id="segment-room-desktop">Camera</button>
            </div>
            <input type="text" id="desktop-single-search"
                   placeholder="Cauta Cladire pe ETTIway" />
            <input type="hidden" id="search-building" />
            <input type="hidden" id="search-room" />
        </div>
        <div id="room-details">
            <h3>Building Details</h3>
            <p class="info-text">Click on a building to view details</p>
        </div>
        <div id="room-list">
            <div class="bottom-sheet-handle"></div>
            <h3>Available Buildings</h3>
            <div id="room-list-container">
                <p class="info-text">Loading buildings...</p>
            </div>
        </div>
    </div>

    <!-- Harta si controalele plutitoare -->
    <div id="map"></div>
    <div id="geo-warning" style="display: none;">
        Atentie: Ai parasit perimetrul campusului ETTI.
    </div>
    <button id="find-me-btn" onclick="toggleLocation()">Find Me</button>
    <button id="recenter-btn" onclick="recenterMap()"
            style="display: none;">Recenter</button>

    <!-- Modale pentru planurile de etaj (continut omis pentru concizie) -->
    <div id="indoor-modal" class="indoor-modal"> ... </div>
    <div id="floor-manager-modal" class="indoor-modal"> ... </div>
    <div id="auth-error-overlay" style="display: none;"> ... </div>

    <script src="https://unpkg.com/leaflet@1.9.4/dist/leaflet.js"></script>
    <script src="js/graph-editor.js"></script>
    <script src="js/utils.js"></script>
    <script src="js/routing.js"></script>
    <script src="js/indoor.js"></script>
    <script src="js/map.js"></script>
    <script src="js/main.js"></script>
</body>
</html>
\end{lstlisting}

\section*{Verificarea rolului pe partea de client}

\begin{lstlisting}[language=JavaScript, caption={Verificarea rolului utilizatorului inainte de activarea modului de editare (defense-in-depth)}, label={lst:authcheck}]
function showAuthError(message) {
    document.getElementById('auth-error-message').textContent = message;
    document.getElementById('auth-error-overlay').style.display = 'flex';
}

async function loadLoggedUser() {
    const isEditMode = new URLSearchParams(window.location.search).get('edit') === '1';
    try {
        const response = await fetch('/api/auth/me', {
            method: 'GET',
            credentials: 'same-origin'
        });

        if (!response.ok) {
            if (isEditMode) {
                showAuthError('Trebuie sa fii logat ca Administrator pentru editare.');
            }
            return;
        }

        const data = await response.json();
        if (data && data.authenticated && data.username) {
            const role = (data.role || 'ROLE_USER').replace('ROLE_', '').toUpperCase();

            // Verificare frontend: doar ADMIN poate intra in modul de editare
            if (isEditMode && role !== 'ADMIN') {
                showAuthError('Contul tau (' + data.username +
                    ') nu are drepturi de Administrator.');
                return;
            }

            const userEl = document.getElementById('logged-user');
            userEl.textContent = 'Logat ca: ' + data.username + ' (' + role + ')';
            userEl.style.display = 'block';
        } else if (isEditMode) {
            showAuthError('Trebuie sa fii logat ca Administrator pentru editare.');
        }
    } catch (error) {
        if (isEditMode) {
            showAuthError('Eroare de conexiune la server.');
        }
    }
}
loadLoggedUser();
\end{lstlisting}
